\chapter{Vektorrally}

Vektorrally er et simpelt spil for to eller flere. Man skal bruge et stykke ternet papir, hvorpå man mere eller mindre tilfældigt tegner den bane hvorpå rallyet skal foregå. Desuden markeres en start/mål-linje, der skal være så bred spillerne kan starte i forskellige gitterpunkter. Når det er en spillers tur, skal hans bil bevæge sig fra et gitterpunkt til et nyt, og man har samme hastighed som man sluttede foregående tur med (altså nul til at starte med). Dog kan man ændre denne hastighed med en ``enhedsvektor", altså en af de otte vektorer (0, $\pm$1), ($\pm$1, $\pm$1), ($\pm$ 1, 0), men man kan naturligvis også undlade at accelerere. Hvis man ikke får bremset eller drejet i tide, og dermed kører af banen, er man ude af spillet. Den spiller som i færrest ture krydser mållinjen vinder. Hvis flere spillere krydser mållinjen i samme antal ture, vinder den som er kørt længst forbi mållinjen.

Et eksempel på en vektorrallyplade og en spillers første par mulige træk er vist nedenfor. Spilleren har en fornuftig retning, og det virker derfor som en god idé at sætte fuld fart på op ad langsiden. Men det er sværere at komme rundt om hjørner end man skulle tro. . .